% CREATED BY MAGNUS GUSTAVER, 2020
\ThesisTitle\\
\ThesisSubtitle\\
\ThesisAuthor\\
\ThesisDepartment\\
Chalmers University of Technology\\
\if\InstitutionLocation G
University of Gothenburg\\
\fi
\setlength{\parskip}{0.5cm}

\thispagestyle{plain}			% Supress header 
\setlength{\parskip}{0pt plus 1.0pt}
\section*{Abstract}
Long-term memory is becoming a central component of large language model agents,
yet existing memory systems often struggle to provide compact and reliable
evidence for temporally grounded questions. In multi-session dialogue histories,
the information needed to answer a query may not be directly named in the query
itself. It can be hidden behind relative temporal expressions, causal
predecessors, or intermediate facts that must first be retrieved and interpreted.
This thesis studies this evidence addressability problem through NanoMem, a
memory framework that treats memory retrieval as iterative evidence synthesis
rather than one-shot context selection or direct answer generation.

NanoMem separates the role of memory from the role of the downstream answer
model. At search time, a Planner rewrites the user query into targeted retrieval
cues, while a Synthesizer compresses retrieved sessions into structured evidence
events. Each event records its source session, the session timestamp, an inferred
event time, and whether the statement is directly supported or inferred. A
verdict signal determines whether the current evidence is sufficient,
indicative, or insufficient; insufficient evidence is fed back to the Planner to
drive another retrieval round. The framework is trained and evaluated with
signals that reward temporal grounding, evidence faithfulness, sufficiency
judgement, and downstream answer quality.

The thesis adapts the completed NanoMem research project and its conference-paper
version into a broader degree-project report. It explains the motivation,
system design, implementation, training signals, benchmark construction, and
experimental evaluation in greater detail than a conference paper permits.
Experiments on long-term memory benchmarks including LoCoMo, LongMemEval, and
ChainMem show that iterative evidence synthesis can improve temporal and
multi-hop memory reasoning while reducing the amount of evidence passed to the
downstream model. The results support the view that agent memory should provide
faithful, source-grounded evidence rather than unstructured retrieved context or
opaque final answers.

% KEYWORDS (MAXIMUM 10 WORDS)
\vfill
Keywords: agent memory, large language models, evidence synthesis, temporal reasoning, retrieval, faithfulness, NanoMem.

\newpage				% Create empty back of side
\thispagestyle{empty}
\mbox{}
